\section{All normalized scales and all centers}\label{sec:scales}

The independent tests of Section~\ref{sec:routing} control a fixed
scale at a fixed stable center.  We first make this estimate simultaneous
over $t\in[1,2]$, using only the grids that determine the local tests.
We then choose the construction parameters and cover the remaining
centers by a small open set.  The finite boundary-representative argument
has precedents in \cite[proof of Theorem~15]{Chlebik2015} and
\cite[Section~3.1]{KolountzakisPapageorgiou2025}.

\begin{lemma}[Conditional control of all normalized scales]
\label{lem:conditional-scales}
Fix $p\in(0,1)$ and the construction parameters $M,d,g,r_0$ of
Sections~\ref{sec:windows} and~\ref{sec:routing}.
Let $x\in\mathcal S$ be a stable center, and let $\xi$ be an atom of
the center-exposure field $\mathcal F_x$ on which the route of $x$
uses a default edge.  If its first default occurs at a vertex $U$ of
height $h$, put $r=r_h$ and
\begin{equation}\label{eq:scale-count}
 D_{M,q}(r)=4+2(M-1)r\bigl(1+2c q^{-2r}\bigr).
\end{equation}
Then
\begin{equation}\label{eq:conditional-scales}
 \Prob\!\left(
  \exists t\in[1,2]\ \forall n\in\mathcal N:\
  x+tq^n\notin B\ \middle|\ \xi
 \right)
 \le D_{M,q}(r)e^{-p(M-1)r/2}.
\end{equation}
\end{lemma}

\begin{proof}
On the atom $\xi$, the vertex $U$ and its outgoing edges
$e_i=(U,i)$, $1\le i<M$, are fixed.  Recall that
$Q_i(z)$ tests the selector on $e_i$ and terminal membership after
routing from $U_i$.  Every table used by this test belongs to $e_i$
or the subtree below it.  By nesting of the grids, for every realization
of the tables, $Q_i(z)$ is determined by the single key
$J_{b_{e_i}^*}(z)$.

For each $i<M$ and $n\in W_{e_i}$, record the values of $t\in[1,2]$
for which $y_n(t)=x+tq^n$ lies on the lattice
$\Gamma_{b_{e_i}^*}$.  The segment traversed by $y_n$ has length
$q^n$, so the number of recorded values is at most
\begin{equation}\label{eq:local-boundaries}
 1+q^nN_{b_{e_i}^*}
 \le 1+2c q^{-2r}.
\end{equation}
Here we used~\eqref{eq:grid-size} and the block-span bound
$b_{e_i}^*-a_{e_i}+1\le2r$ from
Lemma~\ref{lem:window-span}.

Combine these lists, add $1$ and $2$, and retain the distinct values in
increasing order.  Let $\mathcal T_{x,\xi}$ consist of these values
together with one midpoint between each consecutive pair.  There are
$(M-1)r$ pairs $(i,n)$, and hence
\[
 |\mathcal T_{x,\xi}|\le D_{M,q}(r).
\]
This set depends only on $x$, the deterministic windows and grids, and
the vertex $U$ determined by $\xi$; it is fixed before any remaining
table entries are revealed.  For every realization of those entries,
the entire vector
\[
 \bigl(Q_i(y_n(t)):\ i<M,\ n\in W_{e_i}\bigr)
\]
is constant on each open interval between consecutive recorded values.
Every boundary value is represented separately, so the half-open cell
convention causes no loss at a boundary or at an endpoint of $[1,2]$.

If some $t$ misses $B$ at all indices in $\mathcal N$,
Lemma~\ref{lem:routing} implies that all these local tests fail at $t$.
They therefore also fail at some $t'\in\mathcal T_{x,\xi}$.
Each such $t'$ is a fixed real number under the conditioning on $\xi$,
so Lemma~\ref{lem:independent-tests} applies to it.  Taking a union
bound and using~\eqref{eq:test-failure} gives
\[
 \Prob\!\left(
  \exists t\in[1,2]\ \forall n\in\mathcal N:\
  x+tq^n\notin B\ \middle|\ \xi
 \right)
 \le |\mathcal T_{x,\xi}|(1-p/2)^{(M-1)r},
\]
which proves~\eqref{eq:conditional-scales}.
\end{proof}

The bound uses the last grid in each edge and descendant block.
Its resolution is controlled by that edge's window length through
\eqref{eq:window-span}.  Using the finest grid of the whole construction
would discard this control and would not give~\eqref{eq:scale-count}.

\begin{proof}[Proof of Proposition~\ref{prop:periodic}]
Fix $p\in(0,1)$.  We now specify the parameters that were left arbitrary
in the preceding construction.  Choose an integer $M\ge2$ such that
\begin{equation}\label{eq:branching-choice}
 \alpha:=\frac{p(M-1)}2>2\log(1/q).
\end{equation}
Next choose an integer $d\ge1$ satisfying
\[
 (1-2^{1-M})^d<p.
\]
For the resulting tree, let $K$ be its number of edges and choose an
integer $g\ge1$ such that
\[
 4Kc q^{g+1}<p.
\]
Finally, by~\eqref{eq:branching-choice},
\[
 D_{M,q}(r)e^{-\alpha r}
 =O\!\left(r e^{-(\alpha-2\log(1/q))r}\right)
 \longrightarrow0.
\]
Thus there is an integer $r_*\ge1$ such that
\begin{equation}\label{eq:window-length-choice}
 D_{M,q}(r)e^{-p(M-1)r/2}<p
 \qquad\text{for every integer }r\ge r_*.
\end{equation}
Construct the windows with $r_0=r_*$ and the chosen $M,d,g$, and form
the random set $B$ of Section~\ref{sec:routing}.

Fix any stable center $x$.  By~\eqref{eq:no-default}, the probability
that its route has no default is less than $p$.  On each atom of
$\mathcal F_x$ for which a default occurs,
Lemma~\ref{lem:conditional-scales} and~\eqref{eq:window-length-choice}
bound the conditional probability of a miss at some normalized scale
by $p$, since every $r_h\ge r_*$.  Summing over the center-exposure
atoms yields
\begin{equation}\label{eq:stable-miss}
 \Prob\!\left(
  \exists t\in[1,2]\ \forall n\in\mathcal N:\
  x+tq^n\notin B
 \right)\le2p
 \qquad(x\in\mathcal S).
\end{equation}
The probability space is finite, so this event is measurable even
though its definition quantifies over all $t\in[1,2]$.

It remains to turn this estimate at each center into a set that works
at every center. Open-neighborhood repair for null sequences is explicit
in Tom~\cite[Lemma~0.2]{Tom2015}; see also
Chleb\'ik~\cite[Section~3, final remark; Proposition~1(vi), (viii)]{Chlebik2015}.
We implement it with open neighborhoods, which contain sufficiently
late translated points because $tq^n\to0$.
For each table outcome, enlarge the finitely many
selected cells per period to an open periodic set $B^+\supseteq B$
with
\[
 \rho(B^+)\le\rho(B)+p.
\]
For that outcome define the set of centers still missed by the finite
collection of indices:
\begin{equation}\label{eq:exceptional-centers}
 R=\left\{x\in\R:\ \exists t\in[1,2]\ \forall n\in\mathcal N,
                   \ x+tq^n\notin B^+\right\}.
\end{equation}
This is a closed periodic set.  Indeed, on the circle $\R/\Z$ it is
the projection of the closed set
\[
 \left\{(x,t)\in(\R/\Z)\times[1,2]:
       x+tq^n\notin B^+\ \text{for all }n\in\mathcal N\right\}
\]
in a compact product.  In particular, its density is defined.

Since $B\subseteq B^+$, membership of a fixed stable center in $R$
implies the event in~\eqref{eq:stable-miss}.  Lemma~\ref{lem:stability}
and the choice of $g$ give $\rho(\R\setminus\mathcal S)<p$.
Integrating the probability bound over one period therefore gives
\[
 \E\rho(R)
 =\int_0^1\Prob(x\in R)\,dx
 \le2p\rho(\mathcal S)+\rho(\R\setminus\mathcal S)
 \le3p.
\]
Also $\E\rho(B^+)\le2p$ by Lemma~\ref{lem:density}.
Choose a table outcome for which
\[
 \rho(B^+)+\rho(R)\le5p.
\]
Outer regularity on the circle supplies an open periodic set
$V\supseteq R$ with $\rho(V)\le\rho(R)+p$.
Then $H=B^+\cup V$ is open and periodic, and
\[
 \rho(H)\le\rho(B^+)+\rho(R)+p\le6p.
\]

To verify the hitting property, fix $x\in\R$ and $t\in[1,2]$.
If $x\notin R$, the definition of $R$ gives an index $n\in\mathcal N$
with $x+tq^n\in B^+\subseteq H$.
If $x\in R$, then $x\in V$ and $V$ is open, so
$x+tq^n\in V\subseteq H$ for all sufficiently large $n$ because
$tq^n\to0$.  These latter indices need not lie in $\mathcal N$;
the required conclusion allows every $n\ge1$.  Thus every center and
every normalized scale has a hit, proving the proposition.
\end{proof}
